by (i) @@M0@@ (iv), take a surjection @@M1@@. After decomposing @@M2@@ into pieces so that @@M3@@ becomes flat over @@M4@@, the formation of the exact sequence @@M5@@ commutes with passage to the fibers. By (i) @@M6@@ (iii) and (i) @@M7@@ (iv), take @@M8@@ large enough that all the @@M9@@ vanish and that there exist surjections @@M10@@ and @@M11@@. Then the maps @@M12@@ are surjective. Let @@M13@@ be a homomorphism and let @@M14@@ be the composition @@M15@@; then @@M16@@ comes from a homomorphism @@M17@@, and @@M18@@ and @@M19@@ have the same cokernel. Consequently, @@M20@@ is contained in the family of classes of cokernels of homomorphisms of the form @@M21@@.

Again after decomposing @@M22@@ into pieces, one may suppose in addition that @@M23@@ (hence @@M24@@) and all the @@M25@@ (where @@M26@@ is the structural morphism) are flat over @@M27@@, so that the formation of @@M28@@ commutes with arbitrary base change @@M29@@ (EGA III 6.9.9.2). Let @@M30@@; then @@M31@@ canonically carries a ``universal'' exact sequence
@@M32@@
and @@M33@@ bounds @@M34@@.

\sgaStableTarget{sga6.e13.sec.2}\subsection*{2. Several technical lemmas}
\sgaStableTarget{sga6.e13.para.2-1}\textbf{2.1.} Let @@M35@@ be a field and @@M36@@ a proper @@M37@@-scheme. Let @@M38@@ be invertible sheaves on @@M39@@, @@M40@@ the Chern class of @@M41@@ and @@M42@@ a closed subscheme of @@M43@@ of dimension @@M44@@. Recall that we denote by
